% 19-codegen.tex — Chapter 19: Code Generation
\chapter{Code Generation}
\label{chap:19-codegen}
\index{code generation}
\index{C99 backend}

\pnum
The code generator translates the annotated AST into valid C99 source code.  It
is split across \srcref{src/emit/core.h}, \srcref{src/emit/decl.h},
\srcref{src/emit/stmt.h}, \srcref{src/emit/expr.h}, \srcref{src/emit/ctor.h},
\srcref{src/emit/type\_order.h}, and \srcref{src/emit/lain\_header.h}.

% -------------------------------------------------------------------------
\section{Emit context}
\label{sec:gen-ctx}
% -------------------------------------------------------------------------

\pnum
All output is written through a single macro:

\begin{ccode}
static FILE *output_file;
#define EMIT(...)  fprintf(output_file, __VA_ARGS__)
\end{ccode}

\pnum
Indentation is managed by \cname{emit\_indent(depth)}, which writes
\texttt{4 * depth} spaces.  The depth is a plain integer passed through
recursive emit functions.

% -------------------------------------------------------------------------
\section{Output structure}
\label{sec:gen-structure}
% -------------------------------------------------------------------------

\pnum
The generated C file has the following layout:
\begin{enumerate}
\item \textbf{Runtime header}: the content of \srcref{src/emit/lain\_header.h}
  is emitted verbatim.  It defines macro aliases for C library functions
  (\texttt{libc\_printf}, \texttt{libc\_puts}, etc.) and provides the slice
  type primitives.
\item \textbf{\texttt{c\_include} directives}: any \lain{c\_include "foo.h"}
  declarations are emitted as \texttt{\#include "foo.h"}.
\item \textbf{Type forward declarations}: struct and enum types are forward-declared
  via \cname{emit/type\_order.h} in topological dependency order, preventing
  use-before-definition errors in C.
\item \textbf{Struct/enum definitions}: full \texttt{typedef struct} and
  \texttt{typedef enum} bodies.
\item \textbf{Function prototypes}: all functions are forward-declared to allow
  mutual calls.
\item \textbf{Function bodies}: each \cname{DECL\_FUNCTION} and
  \cname{DECL\_PROCEDURE} is emitted with its full body.
\end{enumerate}

% -------------------------------------------------------------------------
\section{Type emission}
\label{sec:gen-types}
\index{type emission}
% -------------------------------------------------------------------------

\pnum
\cname{c\_name\_for\_type(t, out, cap)} converts a Lain \cname{Type *} to a C
type string.  The mapping:

\begin{center}
\begin{tabular}{ll}
\toprule
\textbf{Lain type} & \textbf{C output} \\
\midrule
\lain{int}          & \texttt{int} \\
\lain{u8}           & \texttt{uint8\_t} \\
\lain{i64}          & \texttt{int64\_t} \\
\lain{bool}         & \texttt{bool} \\
\lain{f32}          & \texttt{float} \\
\lain{f64}          & \texttt{double} \\
\lain{u8[5]}        & \texttt{uint8\_t[5]} \\
\lain{u8[:0]}       & \texttt{Slice\_u8\_0} (generated typedef) \\
\lain{mov T}        & Same as \texttt{T} (ownership is erased) \\
\lain{mut T}        & Same as \texttt{T *} (mutable borrow becomes pointer) \\
\bottomrule
\end{tabular}
\end{center}

\pnum
Slice types with sentinels (e.g.\ \lain{u8[:0]}) are represented as C structs
with a \texttt{data} pointer and a \texttt{len} field.  The first use of a slice
type triggers \cname{emit\_slice\_type\_definition} to emit the typedef, which is
then reused for all subsequent uses of the same type.

% -------------------------------------------------------------------------
\section{Enum emission}
\label{sec:gen-enums}
\index{enum emission}
% -------------------------------------------------------------------------

\pnum
Lain enums with payload variants emit:
\begin{enumerate}
\item A \texttt{typedef enum \{...\}} for the tag type.
\item A \texttt{typedef struct \{...\}} with a \texttt{tag} field and a
  \texttt{union} of variant payloads.
\end{enumerate}

\pnum
Constructor functions are generated by \srcref{src/emit/ctor.h} for each variant.
For a variant \lain{Some(value int)}, the constructor is:

\begin{gencode}
Option Option_Some(int value) {
    Option result;
    result.tag = Option_Some_tag;
    result.Some.value = value;
    return result;
}
\end{gencode}

% -------------------------------------------------------------------------
\section{Defer mechanics}
\label{sec:gen-defer}
\index{defer}
% -------------------------------------------------------------------------

\pnum
\lain{defer} statements are implemented using a static stack of up to
\cname{MAX\_DEFERS = 256} statements.  On every return path, the emitter
iterates the defer stack in reverse order and emits each deferred statement
before the \texttt{return}.  Loop-level defer tracking ensures deferred
statements run at loop exit as well as at function exit.

% -------------------------------------------------------------------------
\section{String literal initializers}
\label{sec:gen-strings}
% -------------------------------------------------------------------------

\pnum
String literals initializing typed variables are handled by
\cname{emit\_fixed\_string\_init}.  For a \lain{u8[:0]} target, it emits a
compound literal:

\begin{gencode}
(Slice_u8_0){ .len = 5, .data = (uint8_t[]){ 'h','e','l','l','o', 0 } }
\end{gencode}

\pnum
The sentinel byte is always appended explicitly.  This avoids reliance on
C's implicit NUL-termination of string literals, which may not match the
sentinel value for non-zero sentinels.
